% GENERATED FILE. Edit canonical lessons, then run npm run generate:books.
% canonical-source-hash: fnv1a64:56d096d0a6c32305
% canonical-lessons: GU-C549-phutbol, GU-C549-gol, GU-C549-skor, GU-C549-jhilvu, GU-C549-candrak

\chapter{Goals and Medals}
\label{ch:gu-goals-and-medals}
\hlchaptermodality{549}

\begin{chapteropening}
\textbf{By the end of this chapter:} \emph{I can talk about football, goals, scores and medals.}

Complete the last lesson of chapter 549: I can talk about football, goals, scores and medals.
\end{chapteropening}

\section[\texorpdfstring{\gu{ફૂટબોલ} (phūṭbol)}{ફૂટબોલ (phūṭbol)}]{\gu{ફૂટબોલ} (phūṭbol) — football}

\label{lesson:GU-C549-phutbol}



\begin{quote}
Before the new one: say the Gujarati for a winner, then the Gujarati for to win it for (a team).
\end{quote}

\subsection*{You'll want to know: \gu{ફૂટબોલ}}
\textbf{\gu{ફૂટબોલ}} — \emph{phūṭbol} — \textquotedblleft{}football\textquotedblright{}.

\begin{quote}
\gu{જો} \gu{શનિવારે} \gu{મેદાન} \gu{ખાલી} \gu{હશે} \gu{તો} \gu{અમે} \gu{ફૂટબોલ} \gu{રમીશું}.
\end{quote}

\begin{quote}
\emph{jo śanivāre medān khālī haśe to ame phūṭbol ramīśũ.}
\end{quote}

\begin{quote}
If the ground is free on Saturday, we'll play football.
\end{quote}

A masculine noun.

\begin{cousinweb}
\textbf{\gu{ફૂટબોલ}} (\emph{phūṭbol}) is borrowed from English \emph{football}.
\end{cousinweb}

\begin{grammarlens}[title={What you've built}]
One more word for saying what happens if.
\end{grammarlens}

\subsection*{Your turn}
\begin{itemize}
\raggedright
  \item \emph{Say it:} \emph{phūṭbol}
  \item \emph{Say it:} \emph{phūṭbol}, once more
  \item \emph{Say it:} \emph{jo śanivāre medān khālī haśe to ame phūṭbol ramīśũ.} — If the ground is free on Saturday, we'll play football.
  \item \emph{Recall:} say the Gujarati for a winner, then the Gujarati for to win it for (a team), then say \emph{phūṭbol} again
\end{itemize}

\begin{tcolorbox}[breakable,colback=teal!4,colframe=teal!35!black,title={Before you move on}]
What does \gu{ફૂટબોલ} mean? (\textquotedblleft{}Football\textquotedblright{}.) Say it once more.
\end{tcolorbox}
\section[\texorpdfstring{\gu{ગોલ} (gol)}{ગોલ (gol)}]{\gu{ગોલ} (gol) — a goal (in a game)}

\label{lesson:GU-C549-gol}



\begin{quote}
Before the new one: say the Gujarati for to win it for (a team), then the Gujarati for football.
\end{quote}

\subsection*{You'll want to know: \gu{ગોલ}}
\textbf{\gu{ગોલ}} — \emph{gol} — \textquotedblleft{}a goal (in a game)\textquotedblright{}.

\begin{quote}
\gu{જો} \gu{તું} \gu{એક} \gu{ગોલ} \gu{કરીશ} \gu{તો} \gu{આપણે} \gu{જીતી} \gu{જઈશું}.
\end{quote}

\begin{quote}
\emph{jo tũ ek gol karīś to āpaṇe jītī jaīśũ.}
\end{quote}

\begin{quote}
If you score one goal, we'll win.
\end{quote}

A masculine noun.

\begin{cousinweb}
\textbf{\gu{ગોલ}} (\emph{gol}) is borrowed from English \emph{goal}. One \textquotedblleft{}does\textquotedblright{} a goal: \textbf{\gu{ગોલ} \gu{કરવો}}.
\end{cousinweb}

\begin{grammarlens}[title={What you've built}]
One more word for saying what happens if.
\end{grammarlens}

\subsection*{Your turn}
\begin{itemize}
\raggedright
  \item \emph{Say it:} \emph{gol}
  \item \emph{Say it:} \emph{gol}, once more
  \item \emph{Say it:} \emph{jo tũ ek gol karīś to āpaṇe jītī jaīśũ.} — If you score one goal, we'll win.
  \item \emph{Recall:} say the Gujarati for to win it for (a team), then the Gujarati for football, then say \emph{gol} again
\end{itemize}

\begin{tcolorbox}[breakable,colback=teal!4,colframe=teal!35!black,title={Before you move on}]
What does \gu{ગોલ} mean? (\textquotedblleft{}A goal (in a game)\textquotedblright{}.) Say it once more.
\end{tcolorbox}
\section[\texorpdfstring{\gu{સ્કોર} (skor)}{સ્કોર (skor)}]{\gu{સ્કોર} (skor) — a score}

\label{lesson:GU-C549-skor}



\begin{quote}
Before the new one: say the Gujarati for football, then the Gujarati for a goal (in a game).
\end{quote}

\subsection*{You'll want to know: \gu{સ્કોર}}
\textbf{\gu{સ્કોર}} — \emph{skor} — \textquotedblleft{}a score\textquotedblright{}.

\begin{quote}
\gu{જો} \gu{સ્કોર} \gu{બરાબર} \gu{રહેશે} \gu{તો} \gu{મેચ} \gu{વધારે} \gu{સમય} \gu{ચાલશે}.
\end{quote}

\begin{quote}
\emph{jo skor barābar raheśe to mec vadhāre samay cālśe.}
\end{quote}

\begin{quote}
If the score stays level, the match will go on longer.
\end{quote}

A masculine noun.

\begin{cousinweb}
\textbf{\gu{સ્કોર}} (\emph{skor}) is borrowed from English \emph{score}.
\end{cousinweb}

\begin{grammarlens}[title={What you've built}]
One more word for saying what happens if.
\end{grammarlens}

\subsection*{Your turn}
\begin{itemize}
\raggedright
  \item \emph{Say it:} \emph{skor}
  \item \emph{Say it:} \emph{skor}, once more
  \item \emph{Say it:} \emph{jo skor barābar raheśe to mec vadhāre samay cālśe.} — If the score stays level, the match will go on longer.
  \item \emph{Recall:} say the Gujarati for football, then the Gujarati for a goal (in a game), then say \emph{skor} again
\end{itemize}

\begin{tcolorbox}[breakable,colback=teal!4,colframe=teal!35!black,title={Before you move on}]
What does \gu{સ્કોર} mean? (\textquotedblleft{}A score\textquotedblright{}.) Say it once more.
\end{tcolorbox}
\section[\texorpdfstring{\gu{ઝીલવું} (jhīlvũ)}{ઝીલવું (jhīlvũ)}]{\gu{ઝીલવું} (jhīlvũ) — to catch a ball}

\label{lesson:GU-C549-jhilvu}



\begin{quote}
Before the new one: say the Gujarati for a goal (in a game), then the Gujarati for a score.
\end{quote}

\subsection*{You'll want to know: \gu{ઝીલવું}}
\textbf{\gu{ઝીલવું}} — \emph{jhīlvũ} — \textquotedblleft{}to catch a ball\textquotedblright{}.

\begin{quote}
\gu{જો} \gu{તું} \gu{દડો} \gu{ઝીલત} \gu{તો} \gu{આપણે} \gu{જીતી} \gu{જાત}.
\end{quote}

\begin{quote}
\emph{jo tũ daḍo jhīlat to āpaṇe jītī jāt.}
\end{quote}

\begin{quote}
If you had caught the ball, we would have won.
\end{quote}

Past, with \gu{તેણે} (he or she …): \gu{તેણે} \gu{દડો} \gu{ઝીલ્યો}, \emph{teṇe daḍo jhīlyo}.

\begin{cousinweb}
\textbf{\gu{ઝીલવું}} (\emph{jhīlvũ}) is to catch something thrown or falling. CDIAL lists Marathi \emph{jhelṇẽ}, \textquotedblleft{}to catch\textquotedblright{}, and Hindi \emph{jhelnā}, \textquotedblleft{}to undergo; to take\textquotedblright{}, under a reconstructed \emph{*jhēlati}, \textquotedblleft{}undergoes\textquotedblright{} (CDIAL 5413); it gives no Gujarati form, so its history is not settled.
\end{cousinweb}

\begin{grammarlens}[title={What you've built}]
One more word for saying what happens if.
\end{grammarlens}

\subsection*{Your turn}
\begin{itemize}
\raggedright
  \item \emph{Say it:} \emph{jhīlvũ}
  \item \emph{Say it:} \emph{jhīlvũ}, once more
  \item \emph{Say it:} \emph{jo tũ daḍo jhīlat to āpaṇe jītī jāt.} — If you had caught the ball, we would have won.
  \item \emph{Recall:} say the Gujarati for a goal (in a game), then the Gujarati for a score, then say \emph{jhīlvũ} again
\end{itemize}

\begin{tcolorbox}[breakable,colback=teal!4,colframe=teal!35!black,title={Before you move on}]
What does \gu{ઝીલવું} mean? (\textquotedblleft{}To catch a ball\textquotedblright{}.) Say it once more.
\end{tcolorbox}
\section[\texorpdfstring{\gu{ચંદ્રક} (candrak)}{ચંદ્રક (candrak)}]{\gu{ચંદ્રક} (candrak) — a medal}

\label{lesson:GU-C549-candrak}



\begin{quote}
Before the new one: say the Gujarati for a score, then the Gujarati for to catch a ball.
\end{quote}

\subsection*{You'll want to know: \gu{ચંદ્રક}}
\textbf{\gu{ચંદ્રક}} — \emph{candrak} — \textquotedblleft{}a medal\textquotedblright{}.

\begin{quote}
\gu{જો} \gu{મીરા} \gu{પહેલી} \gu{આવશે} \gu{તો} \gu{તેને} \gu{ચંદ્રક} \gu{મળશે}.
\end{quote}

\begin{quote}
\emph{jo Mīrā pahelī āvśe to tene candrak maḷśe.}
\end{quote}

\begin{quote}
If Mira comes first, she'll get a medal.
\end{quote}

A masculine noun.

\begin{cousinweb}
\textbf{\gu{ચંદ્રક}} (\emph{candrak}) is a tatsama, taken from Sanskrit \emph{candraka}, \textquotedblleft{}the moon; a ring shaped like the moon\textquotedblright{} (Monier-Williams): a round disc. It is built on \textbf{\gu{ચંદ્ર}}, the moon.
\end{cousinweb}

\begin{grammarlens}[title={What you've built}]
One more word for saying what happens if.
\end{grammarlens}

\subsection*{Your turn}
\begin{itemize}
\raggedright
  \item \emph{Say it:} \emph{candrak}
  \item \emph{Say it:} \emph{candrak}, once more
  \item \emph{Say it:} \emph{jo Mīrā pahelī āvśe to tene candrak maḷśe.} — If Mira comes first, she'll get a medal.
  \item \emph{Recall:} say the Gujarati for a score, then the Gujarati for to catch a ball, then say \emph{candrak} again
\end{itemize}

\begin{tcolorbox}[breakable,colback=teal!4,colframe=teal!35!black,title={Before you move on}]
What does \gu{ચંદ્રક} mean? (\textquotedblleft{}A medal\textquotedblright{}.) Say it once more.
\end{tcolorbox}
